\documentclass[a4paper,11pt]{article}
\usepackage{exam-style-341}
\examinehead{341 农业综合知识三 · 数据库分章练习}{第02章\quad 关系模型与关系代数}
\begin{document}

\unitcover{341 农业综合知识三 · 数据库分章练习}{第02章\quad 关系模型与关系代数}{练习（共 100 分）}{约 60 分钟}{关系、元组、属性、域、分量与关系模式；超码、候选码、主码、外码、主属性与非主属性；三类完整性约束；关系代数的五种基本运算与交、连接、除等导出运算；属性集闭包与候选码的求法；关系代数表达式的等价变换。}

\begin{knowmap}
\textbf{关系模型术语：}关系（二维表）· 元组（行）· 属性（列）· 域（属性的取值范围）· 分量（元组中的一个属性值）· 关系模式（对关系的描述）\\[2pt]
\textbf{键体系：}超码 $\supseteq$ 候选码（极小超码）$\rightarrow$ 主码（选定的一个候选码）；外码是本关系中的属性组，参照另一关系的主码；含在任一候选码中的属性为主属性，其余为非主属性\\[2pt]
\textbf{三类完整性：}实体完整性（主属性非空）· 参照完整性（外码取空值或等于被参照元组的主码值）· 用户定义完整性（具体应用的约束）\\[2pt]
\textbf{五种基本运算：}并 $\cup$ · 差 $-$ · 笛卡尔积 $\times$ · 选择 $\sigma$ · 投影 $\pi$；交、连接、除均可由基本运算导出\\[2pt]
\textbf{闭包与候选码：}属性集 $X$ 关于 $F$ 的闭包 $X^{+}$ 是能由 $X$ 经 $F$ 推出的全部属性；若 $X^{+}$ 含全部属性且 $X$ 的任何真子集都不能推出全部属性，则 $X$ 为候选码
\end{knowmap}

\begin{keyterms}
\textbf{关系}：在关系模型中，一个关系就是一张二维表，是元组的集合，元组之间无序、属性之间也无序。\\
\textbf{元组}：关系中的一行，对应现实世界中的一个实体或一条记录。\\
\textbf{属性}：关系中的一列，对应实体的一个特征；属性名在一个关系中必须唯一。\\
\textbf{域}：属性的取值范围，即一组具有相同数据类型的值的集合，又称值域。\\
\textbf{超码}：能够唯一标识关系中一个元组的属性组，可能含有多余的属性。\\
\textbf{候选码}：能够唯一标识一个元组且不含多余属性的属性组，即极小的超码。\\
\textbf{主码}：从关系的多个候选码中选定一个作为主码，用于唯一标识元组；一个关系只有一个主码。\\
\textbf{外码}：本关系中的一个属性组，它不是本关系的主码，但参照另一个关系的主码。\\
\textbf{主属性}：包含在任何一个候选码中的属性；不包含在任何一个候选码中的属性称为非主属性。\\
\textbf{非主属性}：不包含在任何一个候选码中的属性，也称非码属性。\\
\textbf{实体完整性}：若属性（组）$A$ 是基本关系 $R$ 的主属性，则 $A$ 不能取空值。\\
\textbf{参照完整性}：若属性（组）$F$ 是基本关系 $R$ 的外码，则 $R$ 中每个元组在 $F$ 上的值或者取空值，或者等于被参照关系 $S$ 中某个元组的主码值。\\
\textbf{用户定义完整性}：用户针对具体应用环境定义的约束条件，如成绩取值在 0 到 100 之间。\\
\textbf{属性集闭包}：设 $F$ 为关系模式 $R$ 上的函数依赖集，属性集 $X$ 关于 $F$ 的闭包 $X^{+}$ 是由 $X$ 出发、依据 $F$ 中的函数依赖能够推出的全部属性的集合。
\end{keyterms}

\begin{contrast}{超码 / 候选码 / 主码 / 外码（每行依次为：定义·数量关系·要点与示例）}
\citem{超码}{能唯一标识元组的属性组}{可能有多余属性，可有多个}{如（学号，姓名）也能唯一标识学生，但含多余属性}
\citem{候选码}{能唯一标识元组且不含多余属性}{是极小的超码，可有多个}{如学号、身份证号都可唯一标识学生}
\citem{主码}{从候选码中选定的一个}{一个关系只有一个}{数据库设计时指定，如选定学号为主码}
\citem{外码}{参照另一关系主码的属性组}{可以取空值或等于被参照值}{如 SC 中的学号参照 Student 的学号}
\end{contrast}

\begin{contrast}{实体完整性 / 参照完整性 / 用户定义完整性（每行依次为：约束对象·规则内容·违反时的处理）}
\citem{实体完整性}{主属性（主码）}{主属性不能取空值}{拒绝插入或修改}
\citem{参照完整性}{外码}{取空值，或等于被参照关系中某个元组的主码值}{一般拒绝执行，也可按语义级联删除或置空}
\citem{用户定义完整性}{具体应用中的属性}{满足用户规定的取值条件}{拒绝执行}
\end{contrast}

\begin{contrast}{选择 / 投影 / 连接（每行依次为：符号·作用·结果的特点）}
\citem{选择}{$\sigma$}{从关系中选出满足条件的元组（行）}{属性个数不变，元组数可能减少}
\citem{投影}{$\pi$}{从关系中选出指定的属性（列）}{属性个数减少，结果要去掉重复元组}
\citem{连接}{$\bowtie$}{按连接条件把两个关系拼接成一个关系}{属性由两关系属性组合，元组满足连接条件}
\end{contrast}

\begin{contrast}{基本运算 / 导出运算（每行依次为：包含的运算·能否用基本运算表达·作用）}
\citem{基本运算}{并、差、笛卡尔积、选择、投影}{本身不能再被其他运算表达}{关系代数的运算基础}
\citem{导出运算}{交、连接、除等}{都可以用五种基本运算表达}{简化表达，方便书写与理解查询}
\end{contrast}

\qtypename{一、单项选择题}{10 题，每题 2 分，共 20 分}
\q 在关系模型中，二维表中的一行称为（\quad）
\choicesline{A. 属性}{B. 元组}{C. 域}{D. 关系}

\q 能够唯一标识一个元组且不含多余属性的属性组称为（\quad）
\choicesline{A. 超码}{B. 候选码}{C. 主码}{D. 外码}

\q 包含在任何一个候选码中的属性称为（\quad）
\choicesline{A. 主属性}{B. 非主属性}{C. 外码属性}{D. 组合属性}

\q “若属性 $A$ 是基本关系 $R$ 的主属性，则 $A$ 不能取空值”，这条规则属于（\quad）
\choicesline{A. 实体完整性}{B. 参照完整性}{C. 用户定义完整性}{D. 域完整性}

\q 若关系 $R$ 中的属性组 $F$ 是外码，则 $R$ 中每个元组在 $F$ 上的取值（\quad）
\choices{或者为空值，或者等于被参照关系中某个元组的主码值}{必须等于被参照关系中某个元组的主码值}{必须为空值}{可以取任意值}

\q 关系代数中，从关系中选出满足给定条件的元组的运算称为（\quad）
\choicesline{A. 投影}{B. 选择}{C. 连接}{D. 除}

\q 设关系 $R$ 和 $S$ 结构相同，则关系代数表达式 $R-(R-S)$ 等价于（\quad）
\choicesline{A. $R\cup S$}{B. $R\cap S$}{C. $R-S$}{D. $R\times S$}

\q 设关系 $R(A,B)$ 和 $S(B,C)$，则 $R$ 与 $S$ 自然连接后结果的属性个数是（\quad）
\choicesline{A. 2}{B. 3}{C. 4}{D. 5}

\q 下列运算中，不属于关系代数五种基本运算的是（\quad）
\choicesline{A. 并}{B. 差}{C. 交}{D. 投影}

\q 下列关于关系的性质的说法中，正确的是（\quad）
\choices{关系中元组的次序可以任意交换}{关系中允许出现两个完全相同的元组}{关系中的属性可以再分解为更小的数据项}{关系的列次序改变后，关系就变成另一个关系}

\qtypename{二、多项选择题}{5 题，每题 4 分，共 20 分；每题有两个或两个以上正确项}
\q 下列关于关系的性质的说法中，正确的有（\quad）
\choices{每个属性都是不可再分的原子值}{元组的次序可以任意交换}{属性的次序可以任意交换}{不允许出现两个完全相同的元组}

\q 关系数据库的三类完整性约束包括（\quad）
\choicesline{A. 实体完整性}{B. 参照完整性}{C. 用户定义完整性}{D. 数据独立性}

\q 关系代数的五种基本运算包括（\quad）
\choicesline{A. 并}{B. 差}{C. 笛卡尔积}{D. 选择}

\q 关于外码，下列说法正确的有（\quad）
\choices{外码可以是空值}{外码的取值必须等于被参照关系中某个元组的主码值，或者为空值}{外码参照的是另一个关系的主码}{外码必须是本关系的主码}

\q 下列关系代数等价式中，正确的有（\quad）
\choices{$R\cap S=R-(R-S)$}{$R\cup S=S\cup R$}{$R-(R-S)=R\cap S$}{$R-S=S-R$}

\qtypename{三、判断题}{5 题，每题 2 分，共 10 分；正确的画“√”，错误的画“×”并写出改正后的表述}
\q 任何一个二维表都是一个关系。\quad（\quad）\quad 改正：\fixblank

\q 一个关系只能有一个候选码。\quad（\quad）\quad 改正：\fixblank

\q 在参照完整性规则下，外码可以取空值。\quad（\quad）\quad 改正：\fixblank

\q 设关系 $R$ 有 $m$ 个元组，关系 $S$ 有 $n$ 个元组，则 $R\times S$ 有 $m\times n$ 个元组。\quad（\quad）\quad 改正：\fixblank

\q 投影运算的结果中不会出现重复的元组。\quad（\quad）\quad 改正：\fixblank

\qtypename{四、填空题}{5 题，每题 2 分，共 10 分}
\q 关系代数的五种基本运算是并、差、笛卡尔积、\fillblank{} 和 \fillblank{}。

\q 关系代数中，$\sigma$ 表示 \fillblank{} 运算，$\pi$ 表示 \fillblank{} 运算，$\bowtie$ 表示 \fillblank{} 运算。

\q 设关系模式 $R(A,B,C,D)$，函数依赖集 $F=\{A\to B,\ B\to C,\ C\to D\}$，则属性集 $A$ 关于 $F$ 的闭包 $A^{+}=$ \fillblank{}，$R$ 的候选码是 \fillblank{}。

\q 设关系模式 $R(A,B,C)$，函数依赖集 $F=\{AB\to C\}$，则 $R$ 的候选码是 \fillblank{}，主属性是 \fillblank{}，非主属性是 \fillblank{}。

\q 在关系模型中，表中的一行称为 \fillblank{}，表中的一列称为 \fillblank{}，属性的取值范围称为 \fillblank{}。

\qtypename{五、简答题}{2 题，每题 8 分，共 16 分}
\q 试说明超码、候选码、主码和外码的区别与联系，并举例说明。\anslines{2.2cm}

\q 简述关系数据库的三类完整性约束，并分别说明其规则内容和违反时的处理方式。\anslines{2.2cm}

\qtypename{六、综合题}{2 题，每题 12 分，共 24 分}
\q 设某高校学生选课数据库中有三个关系模式（带下划线的属性为主码）：
\begin{center}
\begin{tabular}{@{}ll@{}}
\toprule
关系名 & 结构 \\
\midrule
Student & Student(\underline{Sno}, Sname, Ssex, Sage, Sdept) \\
Course & Course(\underline{Cno}, Cname, Ccredit) \\
SC & SC(\underline{Sno}, \underline{Cno}, Grade) \\
\bottomrule
\end{tabular}
\end{center}

请用关系代数表达式完成下列查询：
\begin{enumerate}[label=（\arabic*）, leftmargin=2.4em, itemsep=1pt]
  \item 查询计算机系（Sdept = 'CS'）年龄小于 20 岁的男学生的姓名。（3 分）
  \item 查询选修了课程号为 'C01' 的课程的学生的姓名。（3 分）
  \item 查询没有选修任何课程的学生的学号。（3 分）
  \item 查询选修了全部课程的学生的学号。（3 分）
\end{enumerate}
\anslines{4.5cm}

\q 设关系模式 $R(A,B,C,D)$，其上的函数依赖集为 $F=\{A\to B,\ B\to A,\ AC\to D\}$。

请完成：
\begin{enumerate}[label=（\arabic*）, leftmargin=2.4em, itemsep=1pt]
  \item 求属性集 $A$、$B$、$C$ 关于 $F$ 的闭包 $A^{+}$、$B^{+}$、$C^{+}$。（3 分）
  \item 求属性集 $AC$、$BC$ 关于 $F$ 的闭包 $AC^{+}$、$BC^{+}$，并由此写出 $R$ 的全部候选码。（4 分）
  \item 指出 $R$ 的主属性和非主属性。（2 分）
  \item 说明 $A$、$B$、$C$ 各自单独为什么不能成为候选码，而 $AC$、$BC$ 为什么是候选码。（3 分）
\end{enumerate}
\anslines{4.5cm}

\end{document}
